% !PRINTED 468
\weberpage{468}

%% [running head, not body text -- commented out 2026-08-07] Page furniture (running
%% head, folio, signature) is excluded from this edition by standing policy.
%% \noindent 468 \hfill Neunzehnter Abschnitt. \hfill \S\,127.

%<<CONT
der eine, der dem Werte $r=0$ oder zwei, die den Werten $r=0,2$ entsprechen, genügen, so folgt, daß $f(\omega)^6$ oder $f(\omega)^{12}$ Klasseninvarianten sind.

Ist außerdem $m$ durch 3 unteilbar, so kann man die $B_r$ alle durch 3 teilbar voraussetzen, und es folgt dann durch Kombination mit dem Satz 1., daß auch
\[
f(\omega)^8f(\omega)^{-6}=f(\omega)^2 \quad\text{oder}\quad f(\omega)^{12}f(\omega)^{-8}=f(\omega)^4
\]
unter den Klasseninvarianten enthalten sind.

Nachdem so unser nächstes Ziel bezeichnet ist, machen wir von dem Ergebnis des \S\,73 Gebrauch, daß zwischen
\[
v=f\!\left(\frac{c+\partial\omega}{a}\right)^{3},\quad u=f(\omega)^3
\]
eine Modulargleichung:
\weq{7}{\Phi_n(u,v)=0}
besteht, wenn
\[
a\partial=n,\quad c\equiv 0\pmod{16}
\]
ist. Die hieraus abgeleitete Gleichung
\weq{8}{\Phi_n(u,u)=0}
ist dann und nur dann befriedigt, wenn
\weq{9}{\dfrac{c+\partial\omega}{a}=\dfrac{\gamma+\delta\omega}{\alpha+\beta\omega},}
worin $\begin{pmatrix}\alpha, & \beta\\ \gamma, & \delta\end{pmatrix}$ eine lineare Transformation der ersten oder zweiten Klasse ist, die [nach \S\,40, (12)] der Bedingung genügt:
\[
\left(\frac{2}{\alpha-\beta}\right)e^{-\frac{\pi i}{8}(\alpha-\beta)(\alpha+\beta+\gamma-\delta)}=1,
\]
oder der damit gleichbedeutenden:
\weq{10}{\alpha\gamma+\beta\delta+2\alpha^2-2\alpha\beta-2\alpha\delta\equiv 0\pmod{16}.}
Aus (9) folgt aber, wenn $\omega$ der Gleichung (2) genügt, wie wir im vorigen Paragraphen gesehen haben,
\weq{11}{
\begin{array}{ll}
n=y^2+mx^2, & \\
\alpha=Bx-y, & n\gamma=-Cx+c(Bx-y),\\
\beta=Ax, & n\delta=-Bx-y+Acx.
\end{array}
}
Nehmen wir zunächst $m$ ungerade an, so können wir $n=m$, also $x=1$, $y=0$ setzen, und da $c$ durch 16 teilbar ist, folgt:
\[
\alpha=B,\quad \beta=A,\quad \gamma\equiv -mC,\quad \delta\equiv -mB\pmod{8},
\]
%CONT>>
